\section{Human Performance in Competitive and Collaborative Human--Machine Teams}

\emph{A short methods excerpt, transcribed into simple TeX from the published article by Bennett, Hedley, Love, Houpt, Brown, and Eidels (2023), DOI 10.1111/tops.12683. Figure references and citations have been omitted. The authors' working TeX source informed the selection.}

\subsection{Task design}

Participants in human--human teams were assigned to control the left or right paddles according to their log-in order, while participants in the human--machine team were randomly allocated starting positions. A player's paddle was always presented in red and their teammate in blue. Paddle movement across the screen was only limited by the vertical edges of the detection-response-task panels, such that the position of one player did not inhibit the movement of the other.

Balls were colored purple in the collaborative and competitive conditions but were separated into red (player) and blue (teammate) in the separate condition. Players were awarded a \textit{hit} when balls rebounded off a paddle and a \textit{miss} when a ball passed the bottom of the screen. Each hit rewarded a player with +1 to their individual score in the separate and competitive conditions or +1 to the combined team score in the collaborative condition.

\subsection{Machine agent design}

Participants in the human--machine task were paired with one of three agent types. The first agent was a reinforcement-learning agent trained using hits as a reward. The second and third machine agents were designed with ideal-observer constructs. The second agent used ball intercept points and times until interception within the task; the third used the same algorithm but was coded to avoid the human player's position in the collaborative condition.

The agents were designed to elicit human-like performance outcomes across workload levels rather than to be optimal task performers.
